\documentclass[11pt,a4paper]{article}

\usepackage[margin=1in]{geometry}
\usepackage{amsmath}
\usepackage{amssymb}
\usepackage{booktabs}
\usepackage{microtype}
\usepackage{algorithm}
\usepackage{algpseudocode}
\usepackage[numbers]{natbib}
\usepackage[hidelinks]{hyperref}

\newcommand{\Q}{\mathbf{Q}}
\newcommand{\Pmat}{\mathbf{P}}
\newcommand{\SubMat}{\mathbf{S}}
\newcommand{\bpi}{\boldsymbol{\pi}}
\newcommand{\tree}{\tau}
\newcommand{\branch}{\mathbf{t}}
\newcommand{\data}{\mathbf{D}}
\newcommand{\states}{\mathcal{A}}

\title{\texttt{phylo}: Technical Document}
\author{M.\ J.\ Wilson}
\date{\today}

\begin{document}
\maketitle

\begin{abstract}
This document states the science behind \texttt{phylo}: the substitution
model, the likelihood and the algorithm that evaluates it, the tree-search
problem and its move sets, the reinforcement-learning formulation layered on
top, and the computational structure that makes any of it tractable. It also
holds the captions for the figures and tables the quality-assurance suite
emits, so that a plot and its description cannot drift apart. Equations and
algorithms follow the formulations in the sources cited throughout;
deviations are stated where they occur.
\end{abstract}

\tableofcontents

\section{Scope and status}

The project's goal is to search the space of phylogenetic tree topologies
with reinforcement learning, scoring candidates by their likelihood under a
model of character substitution. This document develops the theory that
search requires.

The code implements none of it yet. The repository is scaffolding --- build,
test, benchmark, lint, documentation, and audit pipelines around placeholder
functions --- and this document is written ahead of the implementation so
that the implementation has a specification to match. Sections marked
\emph{not yet implemented} describe intent; everything else describes
mathematics that does not depend on our code.

\section{Notation}

Let $\states$ be the finite state space of a single character, with
$k = |\states|$; for nucleotide data $\states = \{A, C, G, T\}$ and $k = 4$.
The data $\data$ are an alignment of $n$ taxa over $m$ sites, so
$\data \in \states^{n \times m}$, with $\data_{\cdot s}$ the observed column
at site $s$.

A phylogeny is a pair $(\tree, \branch)$: a tree topology $\tree$ whose
$n$ leaves carry the taxa, and a vector $\branch$ of non-negative branch
lengths, measured in expected substitutions per site. Internal nodes
represent unobserved ancestors. Table~\ref{tab:symbols} collects the symbols
used throughout.

\begin{table}[htbp]
  \centering
  \begin{tabular}{ll}
    \toprule
    Symbol & Meaning \\
    \midrule
    $\states$, $k$ & character state space and its size \\
    $n$, $m$ & number of taxa, number of sites \\
    $\tree$ & tree topology \\
    $\branch$ & vector of branch lengths \\
    $\Q$ & instantaneous rate matrix, $k \times k$ \\
    $\Pmat(t)$ & transition-probability matrix over a branch of length $t$ \\
    $\bpi$ & root (and, under reversibility, stationary) distribution \\
    $L_u(i)$ & partial likelihood at node $u$ given state $i$ \\
    \bottomrule
  \end{tabular}
  \caption{Symbols used in this document.}
  \label{tab:symbols}
\end{table}

\section{The substitution model}

Character evolution along a branch is a continuous-time Markov chain on
$\states$ with rate matrix $\Q$, whose off-diagonal entries $q_{ij} \ge 0$
give the instantaneous rate of substitution from state $i$ to state $j$ and
whose rows sum to zero,
%
\begin{equation}
  q_{ii} = -\sum_{j \neq i} q_{ij}.
  \label{eq:generator}
\end{equation}
%
Transition probabilities over a branch of length $t$ follow from the matrix
exponential,
%
\begin{equation}
  \Pmat(t) = \exp(\Q t), \qquad
  P_{ij}(t) = \Pr(X_t = j \mid X_0 = i).
  \label{eq:transition}
\end{equation}
%
$\Pmat(t)$ is a stochastic matrix: $P_{ij}(t) \ge 0$ and each row sums to
one. Both properties are invariants the implementation is tested against.

\subsection{Reversibility and the root distribution}

The models used here are time-reversible with respect to a distribution
$\bpi$, meaning they satisfy detailed balance,
%
\begin{equation}
  \pi_i q_{ij} = \pi_j q_{ji} \quad \text{for all } i, j.
  \label{eq:detailed-balance}
\end{equation}
%
Reversibility has two consequences we rely on. First, $\bpi$ is the
stationary distribution of the chain, which makes it the natural choice for
the distribution of states at the root. Second, the likelihood is invariant
to where the root is placed on the tree --- Felsenstein's \emph{pulley
principle} \citep{felsenstein2004} --- so an unrooted topology suffices for
inference, and re-rooting is a test the implementation must pass rather than
a modelling decision.

A general reversible rate matrix factorizes as
%
\begin{equation}
  q_{ij} = s_{ij} \pi_j \quad (i \neq j), \qquad s_{ij} = s_{ji},
  \label{eq:gtr}
\end{equation}
%
which is the general time-reversible (GTR) model; the familiar JC69, K80, and
HKY85 models are constrained special cases \citep{felsenstein2004}. Branch
lengths measure expected substitutions per site, which fixes the scale of
$\Q$ through the normalization
%
\begin{equation}
  -\sum_{i} \pi_i q_{ii} = 1 .
  \label{eq:normalisation}
\end{equation}
%
Without \eqref{eq:normalisation}, $\Q$ and $\branch$ trade off against each
other and neither is identifiable.

\section{The likelihood}

Sites are modelled as independent given $(\tree, \branch, \Q, \bpi)$, so the
log-likelihood is a sum over columns,
%
\begin{equation}
  \log \mathcal{L}(\tree, \branch, \Q, \bpi \mid \data)
    = \sum_{s=1}^{m} \log \Pr(\data_{\cdot s} \mid \tree, \branch, \Q, \bpi).
  \label{eq:site-independence}
\end{equation}
%
The site term marginalizes over the unobserved states of every internal node.
Written directly, that is a sum over $k^{n-1}$ assignments --- intractable
beyond a handful of taxa. Felsenstein's pruning algorithm computes it in time
linear in $n$ by exploiting the tree's conditional independence structure
\citep{felsenstein2004}; it is the sum-product algorithm on a tree
\citep{mackay2003, koller2009}.

\subsection{Pruning}

Define the partial likelihood $L_u(i)$ as the probability of the data at the
leaves below node $u$, given that $u$ is in state $i$. At a leaf with
observed state $a$, $L_u(i) = \mathbb{1}[i = a]$. At an internal node $u$
with children $v$ joined by branches of length $t_v$,
%
\begin{equation}
  L_u(i) = \prod_{v \in \mathrm{children}(u)}
           \left( \sum_{j \in \states} P_{ij}(t_v) \, L_v(j) \right),
  \label{eq:pruning}
\end{equation}
%
and at the root $r$ the site likelihood is
%
\begin{equation}
  \Pr(\data_{\cdot s} \mid \cdot) = \sum_{i \in \states} \pi_i \, L_r(i).
  \label{eq:root}
\end{equation}

\begin{algorithm}[htbp]
\caption{Felsenstein pruning for one site}
\label{alg:pruning}
\begin{algorithmic}[1]
  \Require topology $\tree$, branch lengths $\branch$, rate matrix $\Q$,
           root distribution $\bpi$, observed column $\data_{\cdot s}$
  \Ensure site likelihood $\Pr(\data_{\cdot s} \mid \tree, \branch, \Q, \bpi)$
  \ForAll{nodes $u$ in post-order}
    \If{$u$ is a leaf with observed state $a$}
      \State $L_u(i) \gets \mathbb{1}[i = a]$ for all $i \in \states$
    \Else
      \State $L_u(i) \gets \prod_{v \in \mathrm{children}(u)}
             \sum_{j} P_{ij}(t_v) L_v(j)$ for all $i \in \states$
    \EndIf
  \EndFor
  \State \Return $\sum_{i} \pi_i L_r(i)$
\end{algorithmic}
\end{algorithm}

Algorithm~\ref{alg:pruning} costs $O(n k^2)$ per site against $O(k^{n-1})$
for direct marginalization. The two agree exactly, which makes brute-force
marginalization on small trees the reference implementation the pruning code
is tested against.

Partial likelihoods underflow for realistic $m$ and $n$, so the
implementation will rescale them per node and accumulate the log of the
scaling factors. Rescaling changes the arithmetic but not the result: the
rescaled and direct computations must agree to within floating-point
tolerance on small problems where both run.

\subsection{Rate variation across sites}

Sites evolve at different rates. The standard treatment draws a
site-specific rate multiplier $r$ from a discretized gamma distribution with
mean one, and averages the site likelihood over the $c$ rate categories,
%
\begin{equation}
  \Pr(\data_{\cdot s} \mid \cdot)
    = \frac{1}{c} \sum_{g=1}^{c}
      \Pr(\data_{\cdot s} \mid \tree, r_g \branch, \Q, \bpi),
  \label{eq:gamma-rates}
\end{equation}
%
multiplying the cost of a likelihood evaluation by $c$
\citep{felsenstein2004}. \emph{Not yet implemented.}

\section{Simulation}

Simulation is how the implementation is tested: draw data from known
parameters, then check that each component recovers what the generating model
implies. Generation runs in pre-order, the reverse of pruning.

\begin{algorithm}[htbp]
\caption{Simulating an alignment under $(\tree, \branch, \Q, \bpi)$}
\label{alg:simulate}
\begin{algorithmic}[1]
  \Require topology $\tree$, branch lengths $\branch$, rate matrix $\Q$,
           root distribution $\bpi$, sites $m$, seeded generator
  \Ensure alignment $\data \in \states^{n \times m}$
  \For{$s \gets 1$ to $m$}
    \State draw the root state $x_r \sim \bpi$
    \ForAll{nodes $u$ in pre-order, $u \neq r$ with parent $p$ and branch $t_u$}
      \State draw $x_u \sim \Pmat(t_u)_{x_p, \cdot}$
    \EndFor
    \State record $x_u$ at every leaf $u$ as column $s$
  \EndFor
  \State \Return $\data$
\end{algorithmic}
\end{algorithm}

Every draw uses an explicitly seeded generator, so a simulated dataset is
reproducible from its seed. \emph{Not yet implemented.}

\section{Tree search}

\subsection{The size of the problem}

The number of distinct unrooted binary topologies on $n$ taxa is the double
factorial $(2n - 5)!!$ \citep{rosen2019}, which grows faster than exponentially
(Table~\ref{tab:treecount}). Exhaustive search is impossible past a handful
of taxa, and finding the optimal tree under either the parsimony or the
likelihood criterion is NP-hard \citep{felsenstein2004}. Practical inference
is therefore heuristic: start somewhere, propose rearrangements, keep what
scores better.

\begin{table}[htbp]
  \centering
  \begin{tabular}{rr}
    \toprule
    Taxa $n$ & Unrooted binary topologies $(2n-5)!!$ \\
    \midrule
    10 & $2.0 \times 10^{6}$ \\
    20 & $2.2 \times 10^{20}$ \\
    50 & $2.8 \times 10^{74}$ \\
    \bottomrule
  \end{tabular}
  \caption{Growth of tree space with the number of taxa.}
  \label{tab:treecount}
\end{table}

\subsection{Move sets}

Two rearrangement neighbourhoods define the search graph
\citep{felsenstein2004}:

\begin{description}
  \item[NNI (nearest-neighbour interchange).] Each internal edge of an
    unrooted binary tree separates four subtrees, which can be rejoined in
    three ways; two are new topologies. With $n - 3$ internal edges, the NNI
    neighbourhood contains $2(n-3)$ topologies --- small, cheap, and local.
  \item[SPR (subtree prune and regraft).] Detach a subtree and reattach it at
    another edge. The neighbourhood is quadratic in $n$, so each step costs
    more, but the larger reach escapes local optima that trap NNI.
\end{description}

Since neighbouring topologies differ only near the affected edges, most
partial likelihoods are unchanged between a tree and its neighbour. Caching
and recomputing only the invalidated path is what makes tree search
affordable, and is a design requirement on the likelihood engine rather than
an optimization to add later.

\subsection{Extensions to the move set}
\label{sec:extensions}

Three extensions target the same bottleneck: the number of likelihood
evaluations a search spends. All three are proposals. \emph{Not yet
implemented.}

\paragraph{Bounding whole sets of trees.} Scoring every neighbour exactly
wastes work when most are poor. A cheap \emph{upper} bound on the likelihood
attainable anywhere within a set of topologies lets the search discard the
whole set, in the manner of branch and bound \citep{clrs2022}: when the
bound falls below the best log-likelihood already achieved, no member of the
set can win. Two sources of bounds are worth trying: relaxations of the
pruning recursion \eqref{eq:pruning}, which replace the product over
children by a quantity that is cheaper and provably no smaller; and
compressed summaries of the alignment, where identical site patterns already
collapse and coarser encodings trade tightness for speed
\citep{blahut2003}. Soundness is the whole game --- a bound that can fall
below an attainable likelihood silently discards the optimum --- so each
candidate bound needs a proof, and a test that searches small trees
exhaustively for a counterexample.

\paragraph{Multi-moves.} Composing several NNI or SPR rearrangements into
one action reaches topologies no single move reaches, at the cost of a much
larger action space; in reinforcement-learning terms these are temporally
extended actions \citep{suttonbarto2018}. Rather than fixing how many
rearrangements a compound move contains, draw the composition during
training from a Dirichlet process, so the number of distinct compound moves
the agent maintains grows with the evidence for them instead of being set in
advance \citep{mackay2003}. The measure of success is unchanged: fewer
likelihood evaluations for the same likelihood reached.

\paragraph{A transformer policy over Newick strings.} A topology serializes
to a Newick string, which a transformer can consume directly --- tokenize
the (compressed) string, pretrain on simulated trees, then fine-tune the
policy with the usual policy-gradient machinery
\citep{raschka2024, lapan2020, suttonbarto2018}. One property of the
encoding decides whether this works: Newick is not canonical. The same
topology admits many strings, differing by child order and by the choice of
root for the serialization, so either the encoding is canonicalized or the
model is made invariant to those relabelings. Otherwise the policy learns
distinctions that carry no phylogenetic meaning.

\subsection{Parsimony and likelihood}

The \emph{large parsimony} problem asks for the topology minimizing the
number of state changes needed to explain the data. We keep the search
problem and replace the criterion: candidates are scored by the maximized
log-likelihood of \eqref{eq:site-independence}, with the continuous
parameters $(\branch, \Q, \bpi)$ fitted per topology. Likelihood is the more
principled criterion --- it models the substitution process rather than
counting changes, and it is statistically consistent where parsimony can be
positively misleading \citep{felsenstein2004} --- at the cost of a far more
expensive score.

\section{Gradients}

Fitting $(\branch, \Q, \bpi)$ for a candidate topology is a continuous
optimization, and gradients make it tractable. Differentiating
\eqref{eq:transition} with respect to a branch length gives
%
\begin{equation}
  \frac{\partial \Pmat(t)}{\partial t} = \Q \exp(\Q t) = \Q \, \Pmat(t).
  \label{eq:dpdt}
\end{equation}
%
For a reversible $\Q$ with eigendecomposition $\Q = \mathbf{U} \Lambda
\mathbf{U}^{-1}$, both the exponential and its derivative are diagonal in the
eigenbasis,
%
\begin{equation}
  \Pmat(t) = \mathbf{U} e^{\Lambda t} \mathbf{U}^{-1}, \qquad
  \frac{\partial \Pmat(t)}{\partial t}
    = \mathbf{U} \Lambda e^{\Lambda t} \mathbf{U}^{-1},
  \label{eq:eigen}
\end{equation}
%
so a single decomposition serves every branch. Gradients of the
log-likelihood with respect to all branch lengths follow from reverse-mode
automatic differentiation through the pruning recursion
\eqref{eq:pruning}, at a cost proportional to one forward evaluation
\citep{goodfellow2016, prince2023}. The optimizer consuming them is a
standard constrained quasi-Newton or trust-region method
\citep{nocedal2006}; the parameters carry constraints --- non-negative
branch lengths, a normalized $\bpi$ --- so the geometry of the statistical
manifold, and the natural gradient it induces, is the principled setting
\citep{amari2016}.

Analytic and automatic gradients are checked against central finite
differences,
%
\begin{equation}
  \frac{\partial f}{\partial \theta}
    \approx \frac{f(\theta + h) - f(\theta - h)}{2h},
  \label{eq:finite-difference}
\end{equation}
%
with $h$ chosen to balance truncation against round-off. This check is a
required test for any differentiable component. \emph{Not yet implemented.}

\section{Reinforcement-learning formulation}

Classical heuristics propose moves by a fixed, hand-designed rule. We ask
whether a learned policy proposes better ones, and cast the search as a
Markov decision process $(\mathcal{S}, \mathcal{A}, P, R, \gamma)$
\citep{suttonbarto2018}:

\begin{description}
  \item[State.] The current topology with its fitted continuous parameters,
    together with a summary of the alignment.
  \item[Action.] A move drawn from the NNI or SPR neighbourhood of the
    current topology.
  \item[Transition.] Applying the move and refitting the continuous
    parameters --- deterministic given the action, up to the optimizer.
  \item[Reward.] The improvement in maximized log-likelihood,
    $R = \log \mathcal{L}(\tree') - \log \mathcal{L}(\tree)$.
  \item[Episode.] One search trajectory from a starting tree, terminating on
    a step budget or when no move improves the score.
\end{description}

The natural performance measure is not the likelihood reached but the
likelihood reached \emph{per likelihood evaluation}: every method converges
given unbounded evaluations, so the comparison against a hill-climbing
baseline has to be made at equal budget. \emph{Not yet implemented.}

\section{Computational structure}

Likelihood evaluation dominates the run time: every proposed move costs at
least one evaluation, and search proposes many. Its structure is favourable
--- sites are independent by \eqref{eq:site-independence}, and the per-node
work in \eqref{eq:pruning} is a small dense matrix--vector product --- so the
computation is a batched, data-parallel reduction over a tree.

The project's performance rule follows from that. Where a kernel plausibly
achieves an order of magnitude over the vectorized NumPy reference at the
sizes actually run, it targets the GPU through PyTorch, Triton, or JAX, which
supply autodiff and batching alongside the parallelism; otherwise the work
goes to the Rust backend. The threshold is a measured benchmark, not an
expectation, and every accelerated kernel keeps its reference implementation
as a pinned oracle. Kernel-level considerations --- memory hierarchy,
occupancy, parallel reduction --- follow \citet{pmpp2022}; the CPU-side
counterparts, caching and where time actually goes, follow
\citet{bryant2015}, and the backend itself is Rust \citep{blandy2021}.

Memory is the constraint to watch: partial likelihoods occupy
$O(n \, m \, k \, c)$ floating-point values for $n$ taxa, $m$ sites, $k$
states, and $c$ rate categories, before any caching of unchanged subtrees
across proposals.

\section{Quality assurance}

The test suite is expected to establish scientific validity, not to exercise
syntax: fixtures come from component-wise simulation under known parameters,
expected values come from analytic or brute-force sources, and the invariants
stated above --- stochasticity of $\Pmat(t)$, detailed balance
\eqref{eq:detailed-balance}, agreement of pruning with direct
marginalization, gradients against finite differences \eqref{eq:finite-difference},
re-rooting invariance --- are checked directly.

Part of that suite emits scientific-style figures and tables: parameter
recovery against simulated truth, likelihood surfaces, convergence traces,
and timing comparisons. \textbf{Their captions belong in this section}, added
in the same change that adds the figure, so that a plot and its description
cannot drift apart.

No such figures exist yet --- the quality-assurance framework is unbuilt, and
listing captions for figures that do not exist would be exactly the kind of
placeholder this project's conventions forbid. This section is the place they
go.

\section{Sources}

The texts this project works from, split by what each informs:
\emph{infrastructure} --- how the code is built, structured, and made fast
--- and \emph{application} --- the science it computes. Where an algorithm
we implement appears in one of them, this document follows its formulation
and notation and says where it deviates.

\subsection{Infrastructure}

\begin{description}
  \item[Software craft.] \citet{martin2008} on the structure that keeps a
    change reviewable; \citet{blandy2021} on ownership, lifetimes, and the
    FFI boundary the Rust backend crosses.
  \item[Systems and hardware.] \citet{bryant2015} on the memory hierarchy
    and where CPU time goes; \citet{pmpp2022} on GPU kernels, occupancy,
    and parallel reduction and scan.
  \item[Algorithms and discrete mathematics.] \citet{clrs2022} on data
    structures, traversal, complexity, and branch and bound;
    \citet{rosen2019} on the counting and graph theory behind tree space.
  \item[Numerical optimization.] \citet{nocedal2006} on line search, trust
    region, and quasi-Newton methods for the continuous parameters.
\end{description}

\subsection{Application}

\begin{description}
  \item[Phylogenetics and sequence analysis.] \citet{felsenstein2004} on
    substitution models, pruning, tree search, and the parsimony and
    likelihood criteria; \citet{durbin1998} on probabilistic models of
    sequences; \citet{compeau2015} on the algorithmic view of alignment and
    phylogeny construction; \citet{pachter2005} on the algebraic structure
    of phylogenetic models and their invariants, which offer topology tests
    that avoid likelihood optimization altogether.
  \item[Probabilistic inference and graphical models.] \citet{mackay2003}
    on inference, message passing, and Monte Carlo; \citet{koller2009} on
    factor graphs and the exact-inference machinery that
    \eqref{eq:pruning} instantiates; \citet{frey1998} on belief propagation
    read across inference and coding; \citet{ortega2022} on spectral
    methods for signals indexed by graph vertices.
  \item[Statistical physics and Monte Carlo.] \citet{mezard2009} on
    inference over sparse graphs and phase transitions in hard
    combinatorial problems --- tree search being one; \citet{newman1999} on
    sampling, equilibration, and error estimates from correlated samples.
  \item[Information, coding, and quantum computation.] \citet{blahut2003}
    on algebraic coding, which is the background for the compressed tree
    encodings in Section~\ref{sec:extensions}; \citet{nielsen2010} on
    quantum information and algorithms, background only --- nothing here
    targets quantum hardware.
  \item[Geometry of statistical models.] \citet{amari2016} on the
    statistical manifold, the Fisher metric, and natural gradient.
  \item[Learning and decision making.] \citet{goodfellow2016} and
    \citet{prince2023} on architectures, optimization, and automatic
    differentiation; \citet{suttonbarto2018} on MDPs, value and policy
    methods, and temporally extended actions; \citet{lapan2020} on
    implementation practice; \citet{raschka2024} on transformer internals
    and tokenization, for the policy over Newick strings proposed in
    Section~\ref{sec:extensions}.
\end{description}

\bibliographystyle{plainnat}
\bibliography{references}

\end{document}
